\cleardoublepage
\chapter{IMPLEMENTASI}
\label{chap:implementasi}

\section{Lingkungan Implementasi}
\label{sec:bab5-lingkungan-implementasi}

Implementasi sistem dilakukan pada lingkungan komputasi kelas konsumen dengan memisahkan komponen layanan dan antarmuka pengguna. Pemisahan tersebut digunakan untuk merealisasikan
rancangan sistem pada Bab IV sekaligus menjaga agar komponen inti dapat dikembangkan dan diuji secara terpisah. Lingkungan implementasi mencakup perangkat keras yang digunakan
selama pembangunan dan pengujian serta perangkat lunak yang digunakan untuk menjalankan modul prediksi, pipeline RAG, basis data, dan antarmuka CDSS.

\subsection{Lingkungan Perangkat Keras}
\label{subsec:bab5-lingkungan-perangkat-keras}

Sistem dibangun dan diuji menggunakan laptop ASUS VivoBook M3400QA dengan prosesor AMD Ryzen 5 5600H with Radeon Graphics dan memori sebesar 16 GB. Sistem operasi yang digunakan
adalah Windows 11 Home Single Language 64-bit. Perangkat grafis yang terdeteksi adalah AMD Radeon(TM) Graphics yang terintegrasi. Spesifikasi tersebut merepresentasikan
perangkat kelas konsumen yang digunakan untuk pembangunan dan pengujian sistem.

\begingroup
\setlength{\LTcapwidth}{\textwidth}
\begin{longtable}{ll}
\caption{Spesifikasi perangkat keras lingkungan implementasi.}
\label{tab:tab5.1_spesifikasi_perangkat_keras} \\
\toprule
\textbf{Komponen} & \textbf{Spesifikasi} \\
\midrule
\endfirsthead

\multicolumn{2}{@{}l}{\small\emph{Tabel \thetable{} (lanjutan)}} \\
\toprule
\textbf{Komponen} & \textbf{Spesifikasi} \\
\midrule
\endhead

\bottomrule
\endlastfoot

Perangkat & ASUS VivoBook M3400QA \\

Prosesor & AMD Ryzen 5 5600H with Radeon Graphics \\

Memori & 16 GB RAM \\

Grafis & AMD Radeon(TM) Graphics terintegrasi \\

Sistem operasi & Windows 11 Home Single Language 64-bit \\

Penyimpanan & 512 GB SSD \\

Akselerator & CPU; tidak menggunakan GPU khusus \\
\end{longtable}
\endgroup


Tabel~\ref{tab:tab5.1_spesifikasi_perangkat_keras} merangkum perangkat yang digunakan sebagai lingkungan komputasi utama. Detail identitas mesin seperti nama komputer dan
identitas perangkat tidak dicantumkan karena tidak diperlukan untuk menjelaskan lingkungan implementasi. Implementasi tidak menggunakan akselerator grafis khusus sebagai
komponen komputasi model, sehingga sistem dapat dijalankan pada perangkat kelas konsumen tersebut.

\subsection{Lingkungan Perangkat Lunak}
\label{subsec:bab5-lingkungan-perangkat-lunak}

Lingkungan perangkat lunak disusun sebagai aplikasi web dua lapis. Lapis layanan menggunakan Python 3.11.9 dengan FastAPI dan Uvicorn, sedangkan lapis antarmuka menggunakan
Next.js dan React sebagai aplikasi web progresif. Pengolahan data menggunakan NumPy dan Pandas, model prakiraan menggunakan scikit-learn, penyimpanan potongan dokumen
menggunakan ChromaDB, dan penelusuran leksikal menggunakan Okapi BM25 melalui pustaka \texttt{rank\_bm25}.

Pipeline RAG diorkestrasi menggunakan LangChain, sedangkan penyimpanan data aplikasi menggunakan Supabase. Pustaka \texttt{python-dotenv} digunakan untuk memuat konfigurasi
rahasia dari berkas lingkungan, sementara parameter sistem dipusatkan pada berkas konfigurasi. Pemisahan dependensi layanan dan dependensi eksperimen digunakan agar
komponen yang diperlukan untuk melayani permintaan tidak terbebani oleh pustaka yang hanya diperlukan untuk reproduksi eksperimen.

\begin{table}[h]
    \centering
    \caption{Komponen perangkat lunak lingkungan implementasi sistem.}
    \label{tab:tab5.2_komponen_perangkat_lunak}
    \begin{tabular}{p{0.29\textwidth}p{0.26\textwidth}p{0.35\textwidth}}
        \hline
        \textbf{Peran} & \textbf{Pustaka} & \textbf{Kegunaan} \\
        \hline
        Bahasa dan lingkungan & Python 3.11.9 & Runtime peladen dan skrip \\
        Peladen aplikasi & FastAPI, Uvicorn & API dan peladen ASGI \\
        Antarmuka pengguna & Next.js, React & Aplikasi web progresif \\
        Pengolahan data & NumPy, Pandas & Manipulasi array dan data tabular \\
        Model prakiraan & scikit-learn & Implementasi model gradient boosting \\
        Penyimpanan potongan & ChromaDB & Penyimpanan potongan dokumen dan metadata \\
        Penelusuran leksikal & rank\_bm25 & Penelusuran Okapi BM25 \\
        Orkestrasi RAG & LangChain & Perangkaian pipeline dan koneksi model bahasa \\
        Basis data & Supabase & Penyimpanan data aplikasi \\
        Pemrosesan dokumen & PyPDF2 & Ekstraksi teks dokumen pedoman \\
        Pengujian & pytest & Pengujian otomatis jalur inti \\
        \hline
    \end{tabular}
    \vspace{2pt}
    {\small\itshape Sumber: Dokumentasi lingkungan implementasi sistem.}
\end{table}


Tabel~\ref{tab:tab5.2_komponen_perangkat_lunak} memperlihatkan komponen perangkat lunak yang berperan langsung dalam sistem dan proses pengujian. Pustaka dan komponen yang
hanya digunakan untuk eksperimen reproduksi angka dipisahkan dari dependensi layanan utama. Dengan pengaturan tersebut, lingkungan implementasi dapat menjalankan fungsi
pelayanan tanpa harus memuat seluruh kebutuhan eksperimen.

\subsection{Struktur Sistem dan Konfigurasi}
\label{subsec:bab5-struktur-sistem-konfigurasi}

Perangkat lunak disusun secara modular dengan memisahkan logika inti, layanan API, dan antarmuka pengguna. Logika inti menangani pemrosesan data, prediksi, pembentukan kondisi
klinis, retrieval, dan komponen RAG tanpa bergantung pada antarmuka. Lapis layanan memanggil logika inti tersebut dan menyediakan antarmuka pemrograman untuk digunakan oleh
aplikasi pengguna, sedangkan frontend berkomunikasi dengan lapis layanan melalui API.

Konfigurasi sistem dipisahkan dari implementasi logika sehingga parameter dapat dikelola secara terpusat. Berkas lingkungan digunakan untuk menyimpan informasi sensitif,
sedangkan parameter operasional yang tidak bersifat rahasia ditempatkan pada berkas konfigurasi. Struktur ini digunakan untuk menjaga pemisahan antara kode, konfigurasi,
dan kredensial layanan.

Pada tahap inisialisasi, artefak yang diperlukan untuk prediction dan knowledge retrieval disiapkan sebelum sistem menerima permintaan assessment. Pendekatan tersebut
mengurangi pekerjaan berat yang harus dilakukan pada permintaan pertama dan menjaga waktu respons layanan tetap sesuai dengan karakteristik lingkungan komputasi yang digunakan.
